\section{\texorpdfstring{\texttt{Ring}: bundling an additive \texttt{CommGroup} with multiplication}{Ring: bundling an additive CommGroup with multiplication}}\label{sec:08-rings:03-ring:1}

\subsection{Recall}\label{sec:08-rings:03-ring:2}

Formal definitions cited in this section, gathered here for quick reference (full citations in the \hyperref[sec:root:bibliography:1]{Bibliography}):

\begin{itemize}
\tightlist
\item
  \textbf{Ring.} ``(i) \((R, +)\) is an abelian group, (ii) \(\times\) is associative \ldots{} (iii) the distributive laws hold'' (\cite{DummitFoote2003}, §7.1 ``Basic Definitions and Examples,'' pp.~222--223). A basic consequence: ``\(0a = a0 = 0\) for all \(a \in R\)'' (p.~225, Proposition 1).
\end{itemize}

\begin{lstlisting}
structure Ring (R : Type) where
  addGrp : CommGroup R
  mul : R (*$\rightarrow$*) R (*$\rightarrow$*) R
  one : R
  mul_assoc : (*$\forall$*) a b c : R, mul (mul a b) c = mul a (mul b c)
  one_mul : (*$\forall$*) a : R, mul one a = a
  mul_one : (*$\forall$*) a : R, mul a one = a
  left_distrib : (*$\forall$*) a b c : R, mul a (addGrp.op b c) = addGrp.op (mul a b) (mul a c)
  right_distrib : (*$\forall$*) a b c : R, mul (addGrp.op a b) c = addGrp.op (mul a c) (mul b c)
\end{lstlisting}

Consider each field in turn:

\begin{itemize}
\tightlist
\item
  \texttt{addGrp : CommGroup R} --- the whole additive structure (\(+\), \(0\), unary minus, and commutativity) is a single field, itself a bundled structure. This is the ``structures containing structures'' pattern.
\item
  \texttt{mul}, \texttt{one} --- the multiplicative operation and its identity, \texttt{1}.
\item
  \texttt{mul\_\allowbreak{}assoc}, \texttt{one\_\allowbreak{}mul}, \texttt{mul\_\allowbreak{}one} --- multiplication is associative and has a two-sided identity. Note, however, that \texttt{mul} is \textbf{not} required to be commutative or to have inverses. General rings need neither. (A commutative ring would add a \texttt{mul\_\allowbreak{}comm} field, the same way \texttt{CommGroup} added \texttt{comm} to \texttt{Group}.)
\item
  \texttt{left\_\allowbreak{}distrib}, \texttt{right\_\allowbreak{}distrib} --- multiplication distributes over addition on both sides. Both are needed precisely because \texttt{mul} is not assumed to be commutative.
\end{itemize}

\texttt{addGrp.op} will be written constantly below; notation-free helper abbreviations could be defined later. For now everything is spelled out so each usage is traceable to the definition above.

\begin{mathreading}

\texttt{Ring R} is exactly the textbook definition of a (unital, not-necessarily-commutative) ring, presented as a tuple \[
\Big(\,(R,+,0,-)\in\mathbf{Ab},\ \cdot : R\times R\to R,\ 1\in R,\
\text{proofs of }(\mathrm{R2}),(\mathrm{R3}),(\mathrm{R4})\,\Big).
\] The field \texttt{addGrp} is the \emph{underlying additive abelian group}, so a ring is ``an abelian group \((R,+)\) carrying a compatible monoid structure \((R,\cdot,1)\)'' --- a monoid (a set with an associative operation and identity element, i.e.~a group without inverses). The remaining fields say \((R,\cdot,1)\) is a monoid (\texttt{mul\_\allowbreak{}assoc}, \texttt{one\_\allowbreak{}mul}, \texttt{mul\_\allowbreak{}one}) and that the two operations interact through the two-sided distributive laws --- that is, multiplication is compatible with addition on both sides. Nesting \texttt{addGrp} as a whole substructure mirrors the \hyperref[sec:01-basics:04-terminology:1]{forgetful functor} \(\mathbf{Ring}\to\mathbf{Ab}\) sending a ring to its additive group.

\end{mathreading}

\begin{quote}
Read more: Mathlib's \texttt{Ring} (\texttt{Mathlib.Algebra.Ring.Defs}) sits inside a much larger hierarchy --- \texttt{Semiring}, \texttt{NonUnitalRing}, \texttt{CommRing}, \texttt{DivisionRing}, \texttt{Field} --- each adding or dropping exactly one axiom compared to its neighbors; see \hyperref[sec:13-next-steps:02-moving-to-mathlib:1]{Chapter 13}. For the classical (non-Lean) statement of these axioms and their standard consequences, Dummit \& Foote's \emph{Abstract Algebra} or Aluffi's \emph{Algebra: Chapter 0} (the latter using the same categorical framing this book uses) are standard references.
\end{quote}

\subsection{References}\label{sec:08-rings:03-ring:3}

Full citations in the \hyperref[sec:root:bibliography:1]{Bibliography}. Formal definitions are gathered in Recall, above.

\begin{itemize}
\tightlist
\item
  Dummit and Foote (\cite{DummitFoote2003}), §7.1 ``Basic Definitions and Examples,'' pp.~222--225 --- ring axioms (R1)--(R4) and their consequences.
\item
  Aluffi (\cite{Aluffi2009}) --- \textbf{Gap:} not held in either notebook (see \texttt{NOTEBOOK-\allowbreak{}SOURCE-\allowbreak{}GAPS.md}); this recommendation is offered as further reading, not an independently verified factual claim --- Aluffi's use of forgetful functors and universal properties is publicly documented in the book's own table of contents, not quoted from a verified excerpt.
\end{itemize}
